\documentclass[10pt,a4paper]{article}
\usepackage{nexusS5}
\setnexusheader{Entrée en Troisième — Stage de pré-rentrée}{Évaluation finale — Selim MANSOURI}
\begin{document}
\nexuscover{ÉVALUATION FINALE --- DIAGNOSTIC DE SORTIE}{Durée : 45 minutes}{Selim MANSOURI}{Entrée en Troisième}{Mathématiques --- à traiter individuellement}
\par\noindent\begin{tabularx}{\linewidth}{>{\bfseries}p{22mm} X >{\bfseries}p{16mm} X}
\toprule Nom et prénom & \rule{62mm}{0.32pt} & Date & \rule{34mm}{0.32pt} \\ \bottomrule
\end{tabularx}

\begin{goalbox}
\textbf{Ce document est un diagnostic de sortie, pas un classement.} Il mesure ce qui est
désormais installé, ce qui reste fragile, et il sert à écrire le plan des quatre premières
semaines de la rentrée.
\end{goalbox}

\begin{methodbox}
\textbf{Consignes}
\begin{itemize}
\item Durée : \textbf{45 minutes}. Partie A : environ 10 min. Partie B : environ 20 min. Partie C : environ 15 min.
\item Travailler seul, sans document et sans les cartes de rappel de la séance.
\item Écrire la relation ou la propriété utilisée avant le calcul.
\item Une question peut être laissée de côté puis reprise ; il vaut mieux la laisser vide que d'écrire au hasard.
\item Trois lignes « Certitude » seulement : \conf\ , sur la même échelle qu'en début de stage.
\end{itemize}
\end{methodbox}

\newpage
\phasehead{Partie A --- Automatismes et connaissances essentielles}{environ 10 min}{Questions courtes. Écrire le résultat, sans rédaction longue.}
\begin{enumerate}
\needspace{22mm}
\item[\textbf{A1.}] Calculer $(-15) \div 3 + (-2)\times(-6)$.
\worklines{2}
\par\vspace{1.2mm}
\needspace{22mm}
\item[\textbf{A2.}] Calculer $\dfrac{4}{5}\times\dfrac{15}{8}$ et donner le résultat simplifié.
\worklines{2}
\par\vspace{1.2mm}
\needspace{22mm}
\item[\textbf{A3.}] Réduire $9x - 4 - 5x - 7$.
\worklines{2}
\par\vspace{1.2mm}
\needspace{22mm}
\item[\textbf{A4.}] Résoudre l'équation $4x - 3 = 2x + 9$.
\worklines{2}
\par\vspace{1.2mm}
\needspace{22mm}
\item[\textbf{A5.}] Calculer $(-7)\times(-4) + (-20)\div 5$.
\worklines{2}
\par\vspace{1.2mm}
\needspace{22mm}
\item[\textbf{A6.}] Développer $-5(3x - 2)$.
\worklines{2}
\par\vspace{1.2mm}
\end{enumerate}
\certbox{Certitude sur l'ensemble de la partie A :}
\needspace{68mm}\par\vspace{3mm}
\phasehead{Partie B --- Application et raisonnement}{environ 20 min}{Écrire la relation, la propriété ou la règle avant le calcul, puis contrôler.}
\begin{enumerate}
\needspace{22mm}
\item[\textbf{B1.}] Le triangle $ABC$ est rectangle en $A$. On donne $BC = 20$ cm et $AB = 12$ cm.
\begin{enumerate}
\item Calculer $AC$ en citant le théorème utilisé et en écrivant la relation avant les nombres.
\item Vérifier la vraisemblance du résultat : $AC$ doit-elle être plus petite ou plus grande que $BC$ ? Pourquoi ?
\end{enumerate}
\worklines{5}
\par\vspace{1.2mm}
\needspace{22mm}
\item[\textbf{B2.}] Le triangle $ABC$ est rectangle en $A$. On donne $AB = 5$ cm et $BC = 13$ cm.
\begin{enumerate}
\item Par rapport à l'angle $\widehat{B}$, nommer chacun des côtés $AB$ et $BC$.
\item En déduire $\cos\widehat{B}$, puis la mesure de $\widehat{B}$ arrondie au degré.
\item Si l'on connaissait le côté opposé à $\widehat{B}$ et l'hypoténuse, quel rapport faudrait-il utiliser ?
\end{enumerate}
\worklines{5}
\par\vspace{1.2mm}
\needspace{22mm}
\item[\textbf{B3.}] \begin{enumerate}
\item Un train parcourt $210$ km en $1$ h $30$ min à vitesse constante. Calculer sa vitesse moyenne en km/h.
\item Un élève obtient $14$ (coefficient $2$) et $9$ (coefficient $3$). Calculer sa moyenne pondérée, puis vérifier
qu'elle est comprise entre $9$ et $14$.
\end{enumerate}
\worklines{5}
\par\vspace{1.2mm}
\needspace{22mm}
\item[\textbf{B4.}] \begin{enumerate}
\item Sept cahiers identiques coûtent $21$ TND. Calculer le prix de onze cahiers en passant explicitement par le prix unitaire.
\item Le triangle $ABC$ est rectangle en $A$, avec $AB = 12$ cm et $AC = 16$ cm. Calculer $BC$.
\end{enumerate}
\worklines{6}
\par\vspace{1.2mm}
\end{enumerate}
\certbox{Certitude sur la question B2 :}
\needspace{68mm}\par\vspace{3mm}
\phasehead{Partie C --- Transfert et synthèse}{environ 15 min}{Reconnaître la notion utile, mobiliser plusieurs acquis, conclure par une phrase.}
\begin{enumerate}
\needspace{22mm}
\item[\textbf{C1.}] \textbf{La rampe d'accès.} Une rampe d'accès s'appuie sur une longueur au sol de $4{,}8$ m
et atteint une hauteur de $1{,}4$ m. Le triangle formé est rectangle.
\begin{enumerate}
\item Calculer la longueur de la rampe en citant le théorème utilisé.
\item Calculer la mesure de l'angle formé par la rampe et le sol, arrondie au degré.
\item La réglementation impose une pente d'au plus $5\,\%$, la pente étant le quotient de la hauteur par la
longueur au sol. Calculer la pente en pourcentage, puis conclure.
\item Le coût de fabrication est donné par $C(x) = 85x + 120$ TND, où $x$ est la longueur de la rampe en mètres.
Calculer $C$ pour la rampe étudiée et expliquer ce que représente le nombre $120$.
\end{enumerate}
\worklines{12}
\par\vspace{1.2mm}
\needspace{22mm}
\item[\textbf{C2.}] \begin{enumerate}
\item Développer et réduire $(x + 2)(x + 5)$.
\item Contrôler le résultat pour $x = 1$.
\end{enumerate}
\worklines{5}
\par\vspace{1.2mm}
\end{enumerate}
\certbox{Certitude sur la question C1 :}
\brouillon{\dimexpr\textheight/3\relax}
\newpage
\sectiontitle{Après l'évaluation --- à remplir en deux minutes}
\begin{methodbox}Ces trois lignes ne sont pas notées. Elles servent à construire le plan des quatre premières semaines de la rentrée.\end{methodbox}
\par\noindent\begin{tabularx}{\linewidth}{>{\bfseries}p{62mm} X}
\toprule
La question que j'ai trouvée la plus sûre & \rule{85mm}{0.32pt} \\
La question sur laquelle j'ai le plus hésité & \rule{85mm}{0.32pt} \\
Une notion que je veux revoir dès la première semaine & \rule{85mm}{0.32pt} \\
\bottomrule\end{tabularx}
\vfill
\begin{graybox}\small Ce document est un diagnostic de sortie. Il sera analysé compétence par compétence, puis comparé au positionnement passé en début de stage lorsque cette comparaison est légitime.\end{graybox}
\end{document}
